Conventions as in Theorem \claimref{mass.exact-join}: the chart's complex unit, written $i$, is $I_{6}$, and the bar is the conjugation that reverses it; $Z = (\omega, \pi)$ with $\omega \in \Som$, $\pi \in \Spi$; $\hp{\pi}{\omega} = \sum_{a}\overline{\pi_{a}}\,\omega^{a}$; the twistor form is $\bar{Z}\cdot W = \hp{\pi_{Z}}{\omega_{W}} + \overline{\hp{\pi_{W}}{\omega_{Z}}}$; the moment of the composite is $M_{A} = i\sum_{i}\bar{Z}^{i}\cdot AZ^{i}$. Write $\mathbf{Z}$ for the $4\times N$ matrix with columns $Z^{i}$ and $S$ for the matrix of the twistor form, $\bar{Z}\cdot W = Z^{\dagger}SW$; $S$ exchanges the two halves, so it has eigenvalues $+1, +1, -1, -1$, and $h = \mathbf{Z}^{\dagger}S\mathbf{Z}$. Write $\Omega_{\lambda}$ for the matrix with $Z^{\mathsf T}\Omega_{\lambda}W = \br{\pi_{Z}\pi_{W}} + \lambda[\omega_{Z}\omega_{W}]$, that is $\Omega_{\lambda} = \operatorname{diag}(\lambda\varepsilon, \varepsilon)$ in the blocks $(\Som, \Spi)$, so $B^{il}_{\Lambda} = (Z^{i})^{\mathsf T}\Omega_{\lambda}Z^{l}$ (Theorem \claimref{mass.exact-join}(1), whose derivation is per pair).

(1) \emph{Slots.} The decomposition $Q(P) = -2\lambda d^{2} + Q(p) + 2\lambda\,p\cdot k + \lambda^{2}\,Q(k)$ in the proof of Theorem \claimref{mass.exact-join}(2) uses only the linearity of $U\wedge I \mapsto M_{U\wedge I}$ and the ambient pairing in the chart ($n_{\infty}, n_{0}$ null with $\ip{n_{\infty}}{n_{0}} = 1$, the $x$-slot orthogonal to both with $Q(x) = -2\det x$, Fact \claimref{mass.imported-flat-model-twistor-chart-infinity}, so that the dual metric on forms is $f\cdot g = -\varepsilon_{ab}\varepsilon_{a'b'}f_{aa'}g_{bb'}$); it does not refer to the number of constituents and holds verbatim. The moment is a sum over constituents, so $p = \sum_{i}p^{i}$, $k = \sum_{i}k^{i}$, $d = \sum_{i}d^{i}$, where by Lemma \claimref{mass.chart-form-spin-isomorphism}(2), (3), (1), exactly as in that proof, $p^{i}_{aa'} = i\,\overline{\pi^{i}_{a}}\,(\varepsilon\pi^{i})_{a'}$, $k^{i}_{aa'} = -i\,(\varepsilon\omega^{i})_{a}\,\overline{\omega^{i}_{a'}}$ and $d^{i} = \mathrm{Im}\,t_{ii}$. The three pairings computed there, $p^{i}\cdot p^{l} = |\br{\pi^{i}\pi^{l}}|^{2}$, $k^{i}\cdot k^{l} = |[\omega^{i}\omega^{l}]|^{2}$ and $p^{i}\cdot k^{l} = |t_{il}|^{2}$, use only the $2\times 2$ identities $\varepsilon^{\mathsf T}\varepsilon = 1$, $\varepsilon^{2} = -1$, $\det\varepsilon = 1$ and hold for any pair $(i, l)$. By bilinearity, $Q(p) = \sum_{i,l}|\br{\pi^{i}\pi^{l}}|^{2} = 2\sum_{i<l}|\br{\pi^{i}\pi^{l}}|^{2}$ (the bracket is antisymmetric, so the diagonal vanishes and the off-diagonal terms come in equal pairs), likewise $Q(k) = 2\sum_{i<l}|[\omega^{i}\omega^{l}]|^{2}$, $p\cdot k = \sum_{i,l}|t_{il}|^{2} = \tr(tt^{\dagger})$, and $d = \sum_{i}\mathrm{Im}\,t_{ii} = \mathrm{Im}\tr t$. This is (1).

(2) \emph{The comparison.} Expanding,
\[
2\sum_{i<l}|B^{il}_{\Lambda}|^{2} = Q(p) + \lambda^{2}\,Q(k) + 4\lambda\sum_{i<l}\mathrm{Re}\bigl(\br{\pi^{i}\pi^{l}}\,\overline{[\omega^{i}\omega^{l}]}\bigr).
\]
For $i < l$ let $t^{(il)}$ be the principal $2\times 2$ submatrix of $t$ on rows and columns $i, l$. It is the product of the matrix with rows $\overline{\pi^{i}}, \overline{\pi^{l}}$ and the matrix with columns $\omega^{i}, \omega^{l}$, so $\det t^{(il)} = \overline{\br{\pi^{i}\pi^{l}}}\,[\omega^{i}\omega^{l}]$, as in Theorem \claimref{mass.exact-join} for the pair $(1,2)$, and $\mathrm{Re}\bigl(\br{\pi^{i}\pi^{l}}\overline{[\omega^{i}\omega^{l}]}\bigr) = \mathrm{Re}\det t^{(il)}$. The sum of the principal $2\times 2$ minors of a square matrix is its second elementary symmetric function, $\sum_{i<l}\det t^{(il)} = \tfrac12\bigl((\tr t)^{2} - \tr(t^{2})\bigr)$. Hence, with (1),
\[
Q(P) - 2\sum_{i<l}|B^{il}_{\Lambda}|^{2} = 2\lambda\Bigl(\tr(tt^{\dagger}) - (\mathrm{Im}\tr t)^{2} - \mathrm{Re}(\tr t)^{2} + \mathrm{Re}\tr(t^{2})\Bigr) = 2\lambda\Bigl(\tr(tt^{\dagger}) + \mathrm{Re}\tr(t^{2}) - (\mathrm{Re}\tr t)^{2}\Bigr),
\]
using $\mathrm{Re}(z^{2}) = (\mathrm{Re}\,z)^{2} - (\mathrm{Im}\,z)^{2}$. From $h = t + t^{\dagger}$: $\tr(h^{2}) = \tr(t^{2}) + \tr((t^{\dagger})^{2}) + 2\tr(tt^{\dagger}) = 2\,\mathrm{Re}\tr(t^{2}) + 2\tr(tt^{\dagger})$ and $\tr h = 2\,\mathrm{Re}\tr t$, so the bracket is $\tfrac12\tr(h^{2}) - \tfrac14(\tr h)^{2} = C_{N}$. Finally $2\lambda = Q(I_{\Lambda}) = \Lambda/3$ by Lemma \claimref{mass.de-sitter-scale-frozen-point}(3). With $\tr(h^{2}) = \tr(h_{0}^{2}) + (\tr h)^{2}/N$ the second form of $C_{N}$ follows, and at $N = 2$, $C_{2} = \tfrac14(h_{11} - h_{22})^{2} + |h_{12}|^{2} = C$.

(3) \emph{Invariance.} Under $\mathbf{Z} \mapsto \mathbf{Z}U$: $h \mapsto U^{\dagger}hU$ and $B_{\Lambda} = \mathbf{Z}^{\mathsf T}\Omega_{\lambda}\mathbf{Z} \mapsto U^{\mathsf T}B_{\Lambda}U$; the moment $M_{A} = i\tr(\mathbf{Z}^{\dagger}SA\mathbf{Z})$ is unchanged because $UU^{\dagger} = 1$, so the grip and $m^{2}$ are invariant; $\tr(h^{2})$, $\tr h$ and $\tr(B_{\Lambda}B_{\Lambda}^{\dagger}) = 2\sum_{i<l}|B^{il}_{\Lambda}|^{2}$ are invariant under unitary conjugation and congruence. For $N > 2$ the coefficient $-(N-2)/(4N)$ of $(\tr h)^{2}$ in the second form is nonzero, so $C_{N}$ is not a multiple of the $\su(N)$ Casimir $\tfrac12\tr(h_{0}^{2})$.

(4) \emph{Positivity.} If $h$ is positive definite on a subspace $V \subset \C^{N}$, then $\mathbf{Z}$ is injective on $V$ (as $\mathbf{Z}v = 0$ gives $v^{\dagger}hv = 0$) and $S$ is positive definite on $\mathbf{Z}V$, so $\dim V \le 2$; likewise for negative definite. So the eigenvalues of $h$ are $a_{1} \ge a_{2} \ge 0$, $-b_{1}, -b_{2}$ with $b_{1}, b_{2} \ge 0$, and zeros. With $A = a_{1} + a_{2}$, $B = b_{1} + b_{2}$,
\[
4C_{N} = 2\sum\mu^{2} - \Bigl(\sum\mu\Bigr)^{2} = \bigl(2a_{1}^{2} + 2a_{2}^{2} - A^{2}\bigr) + \bigl(2b_{1}^{2} + 2b_{2}^{2} - B^{2}\bigr) + 2AB = (a_{1} - a_{2})^{2} + (b_{1} - b_{2})^{2} + 2AB \ge 0,
\]
and with (2), $m^{2} \ge 0$ for $\Lambda \ge 0$. Equality $C_{N} = 0$ holds iff $a_{1} = a_{2}$, $b_{1} = b_{2}$ and $AB = 0$, that is iff $h = 0$ or $h = cP$ with $P$ an orthogonal projector of rank two and $c \ne 0$ real. For the Hermitian matrix $1_{N}$, $C_{N} = N/2 - N^{2}/4 < 0$ when $N \ge 3$; it has $N$ positive eigenvalues and is not a Gram matrix of conformal spinors.

(5) \emph{The massless locus.} Define the antilinear map $J Z := S\Omega_{\lambda}^{\mathsf T}\bar{Z}$, with $\bar{Z}$ the componentwise conjugate. Since $S$ and $\Omega_{\lambda}$ are real, $S^{2} = 1$ and $\Omega_{\lambda}^{\mathsf T} = -\Omega_{\lambda}$: (a) $(JZ)^{\dagger}SW = Z^{\mathsf T}\Omega_{\lambda}S^{2}W = Z^{\mathsf T}\Omega_{\lambda}W$; (b) $J^{2}Z = S\Omega_{\lambda}^{\mathsf T}S\Omega_{\lambda}^{\mathsf T}Z$ and $S\Omega_{\lambda}^{\mathsf T}S = \operatorname{diag}(\varepsilon^{\mathsf T}, \lambda\varepsilon^{\mathsf T})$, so $J^{2} = \operatorname{diag}(\lambda(\varepsilon^{\mathsf T})^{2}, \lambda(\varepsilon^{\mathsf T})^{2}) = -\lambda$; (c) applying (a), the antisymmetry of $\Omega_{\lambda}$, (a) again with $JW$ in place of $Z$, and (b): $(JZ)^{\dagger}S(JW) = Z^{\mathsf T}\Omega_{\lambda}(JW) = -(JW)^{\mathsf T}\Omega_{\lambda}Z = -(J^{2}W)^{\dagger}SZ = \lambda\,W^{\dagger}SZ = \lambda\,\overline{Z^{\dagger}SW}$. Let $\Lambda > 0$ and $m^{2} = 0$. By (2) and (4) all $B^{il}_{\Lambda}$ vanish and $C_{N} = 0$. Suppose $h = cP$ with $c \ne 0$, and let $v_{1}, v_{2}$ be orthonormal in the range of $P$. The vectors $\mathbf{Z}v_{1}, \mathbf{Z}v_{2}$ have Gram matrix $c\,1_{2}$, so they span a plane $L$ on which the twistor form is definite with the sign of $c$. The form $Z^{\mathsf T}\Omega_{\lambda}W$ vanishes on the span of the $Z^{i}$, hence on $L$, so by (a) $JL$ is orthogonal to $L$ for the twistor form. By (b) $J$ is injective, so $JL$ is a plane (a complex subspace, $J$ being antilinear), and by (c) the twistor form is definite on $JL$ with the sign of $\lambda c$, the sign of $c$. Then $L \oplus JL$ is a four-dimensional subspace on which the twistor form is definite, which contradicts its signature $(2,2)$. Hence $h = 0$. Conversely, if all $B^{il}_{\Lambda}$ vanish and $h = 0$ then $m^{2} = 0$ by (2). At $N = 2$ this is Theorem \claimref{mass.exact-join}(5) with the multiple of the identity forced to be zero; the examples $Z^{i} = \alpha_{i}e_{1} + \beta_{i}e_{4}$ have every bracket and every entry of $h$ zero, for any $N$.

(6) \emph{The de Sitter Casimir.} By Proposition \claimref{mass.casimir-form-join}(1) and (2), for any number of constituents $Q(P) = Q(I_{\Lambda})\bigl(C_{2}(\Lambda^{2}\CS) - C_{2}(I_{\Lambda}^{\perp})\bigr)$ and $C_{2}(\Lambda^{2}\CS) = \tr(h^{2}) - \tfrac14(\tr h)^{2}$. With (2) and $Q(I_{\Lambda}) = 2\lambda \ne 0$,
\[
C_{2}(I_{\Lambda}^{\perp}) = \tr(h^{2}) - \tfrac14(\tr h)^{2} - \frac{1}{\lambda}\sum_{i<l}|B^{il}_{\Lambda}|^{2} - C_{N} = \tfrac12\tr(h^{2}) - \frac{6}{\Lambda}\sum_{i<l}|B^{il}_{\Lambda}|^{2};
\]
at $N = 2$, $\tfrac12\tr(h^{2}) = \tfrac14(\tr h)^{2} + C$ and this is Proposition \claimref{mass.casimir-form-join}(3).

The script assembles the grip three ways for $N = 2, \dots, 5$ in exact Gaussian-rational arithmetic (from the definition with the chart's matrices and a generic inverse of the ambient Gram matrix; from $\sum|B^{il}_{\Lambda}|^{2}$ with the brackets as determinants $\varepsilon(I_{\Lambda}\wedge Z^{i}\wedge Z^{l})$ and $C_{N}$ from $h$; from the component formulas), and confirms (1) and (2) at four values of $\lambda$, (3) under exact Cayley unitaries, (4) on the samples together with the eigenvalue identity and $C_{N}(1_{N}) < 0$, the three properties of $J$, the massless examples, and (6).
